\section{Keep the signed residual, not just its magnitude}
SAD discards the sign of every difference because it needs a comparison score.
A transform needs the differences themselves. If a current sample is 90 and its
prediction is 100, its residual is $90-100=-10$. A residual of $+10$ would tell
the decoder to increase the prediction instead. Keep $e=C-P$ in signed storage;
use $|e|$ only when calculating a score.

Eight-bit input samples lie between zero and 255. A prediction formed from those
samples remains in that range, and a residual lies between $-255$ and $255$.
A signed byte cannot hold that whole interval. A signed halfword can. Store
sixty-four halfwords for an eight-by-eight residual block, in row-major order,
without clipping or taking magnitudes. Later transform and quantization stages
will change its representation. They do not justify discarding its sign now.

Before quantization, $P+e=C$ exactly. Check this identity at every sample, not
only through a block checksum. Different wrong residuals can have the same total
magnitude. This exact reconstruction is a diagnostic for the current arithmetic;
it is not a promise that the eventual compressed picture will be lossless.

\section{Use the same interpolation as the decoder}
Let $A$ be the integer-position reference sample at the supplied pointer, $B$
its right neighbour, $C$ the sample below it and $D$ the sample below-right.
The two phase bits $p_x,p_y$ select integer or half-sample positions. For this
baseline profile, use
\[
\begin{array}{ll}
P_{0,0}=A,&P_{1,0}=\lfloor(A+B+1)/2\rfloor,\\[3pt]
P_{0,1}=\lfloor(A+C+1)/2\rfloor,&
P_{1,1}=\lfloor(A+B+C+D+2)/4\rfloor.
\end{array}
\]
These are the default rounding rules in H.263 clause 6.1.2, with the profile's
rounding control equal to zero. Each output pixel applies the same selected
formula at its own reference origin. Phase one means half of a sample interval,
not an extra whole-sample displacement.

Sum the four diagonal neighbours before rounding. Two rounded pair averages,
followed by another rounded average, can produce a different answer. With
$A=B=C=0$ and $D=1$, the prescribed diagonal value is
$\lfloor(1+2)/4\rfloor=0$. Averaging the two rows with upward rounding gives zero
and one, and averaging those results with upward rounding gives one.

\begin{figure}[htbp]
\centering
\begin{tikzpicture}[every node/.style={font=\small},alt={Four reference samples surround a half-sample position: the top-left, top-right and bottom-left samples are zero; the bottom-right sample is one. Summing once and adding two before division by four predicts zero at the centre. Rounding the two rows separately and then averaging would incorrectly predict one.}]
\draw[gray] (0,0) rectangle (4,2);
\draw[dashed,gray] (0,1)--(4,1) (2,0)--(2,2);
\node[fill=white] at (0,2) {$A=0$};
\node[fill=white] at (4,2) {$B=0$};
\node[fill=white] at (0,0) {$C=0$};
\node[fill=white] at (4,0) {$D=1$};
\node[draw,fill=white,align=center] at (2,1) {half-sample position\\$P=0$};
\end{tikzpicture}
\caption{A one-sample difference exposes double rounding in diagonal prediction.}
\end{figure}

The four-neighbour sum is at most 1020; adding its rounding bias gives at most
1022. Unsigned halfword lanes therefore have enough range for the entire sum.
Byte addition would lose information before the division. For eight adjacent
outputs, gather the eight unsigned neighbours into halfword lanes of each input
vector. With current samples in \texttt{\$w0} and neighbours $A,B,C,D$ in
\texttt{\$w1} through \texttt{\$w4}, a diagonal arithmetic segment is:
\begin{minted}{mars-idea}
addv.h $w5, $w1, $w2
addv.h $w6, $w3, $w4
addv.h $w5, $w5, $w6
srlri.h $w5, $w5, 2
subv.h $w6, $w0, $w5
\end{minted}
The rounded logical shift adds the single required bias before shifting by two.
The final halfword subtraction keeps signed residual bit patterns, including
negative values. A separate magnitude calculation can accumulate SAD without
changing those residuals.

\section{Derive chroma origins with signed arithmetic}
The search result is expressed in full luma pixels. In the selected profile,
its chroma displacement is half that value. A full-pixel luma displacement of
one therefore needs a half-sample chroma phase. The standard's quarter-position
modification does not introduce another case here: the chosen full-pixel luma
search produces only integer or half-sample chroma positions.

For either signed component $d$, split $d/2$ into an integer origin shift $q$ and
a phase $p$ with
\[
q=\lfloor d/2\rfloor,\qquad p=d-2q.
\]
The phase is zero or one, even for negative $d$. Arithmetic right shift by one
forms $q$ for a full-width signed register; its low bit supplies $p$. Signed
integer division that truncates toward zero does not give this decomposition
for a negative odd value.
\begin{center}
\tagpdfsetup{table/header-rows={1}}
\begin{tabular}{rrrr}
\textbf{Luma $d$} & \textbf{Chroma $d/2$} & \textbf{Origin $q$} & \textbf{Phase $p$}\\
$-2$ & $-1$ & $-1$ & 0\\
$-1$ & $-0.5$ & $-1$ & 1\\
0 & 0 & 0 & 0\\
1 & 0.5 & 0 & 1\\
2 & 1 & 1 & 0
\end{tabular}
\end{center}
At a macroblock whose luma origin is $(16,32)$, the chroma origin is $(8,16)$.
Luma displacement $(-1,-1)$ selects chroma origin $(7,15)$ with phases $(1,1)$.
The top-left prediction then uses the four surrounding chroma samples. Using
origin $(8,16)$ with the same phases would instead select the wrong positive
half-sample displacement.

\section{Keep six transform blocks distinct}
A sixteen-by-sixteen luma macroblock contains four eight-by-eight luma blocks.
Its two eight-by-eight chroma blocks bring the total to six. Use the following
order when preparing the later coefficient stage:
\begin{center}
\tagpdfsetup{table/header-rows={1}}
\begin{tabular}{rll}
\textbf{Block} & \textbf{Plane} & \textbf{Offset within that plane's macroblock}\\
0 & Y & $(0,0)$\\
1 & Y & $(8,0)$\\
2 & Y & $(0,8)$\\
3 & Y & $(8,8)$\\
4 & Cb & $(0,0)$\\
5 & Cr & $(0,0)$
\end{tabular}
\end{center}
For the four luma blocks, apply the selected full-pixel displacement and use
phases zero. Add the block's local offset to both current and reference origins.
For Cb and Cr, derive $q_x,q_y,p_x,p_y$ once from the same luma vector, then use
the appropriate chroma plane and stride. The two planes share geometry, not sample
values. Do not reuse the Cb pointer for Cr.

A fractional phase needs an extra reference neighbour: the active footprint is
$(8+p_x)$ columns by $(8+p_y)$ rows. This is nine by nine for diagonal prediction.
The original sixteen-by-sixteen luma search bounds keep the corresponding chroma
neighbours within the picture for aligned macroblock origins in this restricted
full-pixel search. A general caller of the following kernel must still supply
its stated mapped footprint. The kernel does not infer picture geometry from a
pointer or manufacture missing border samples.
